\section{实验设计}
\label{sec:experiments}

\subsection{组织原则与可追溯性}
本文全部数值结果由脚本生成、可复算：每次实验以固定 \texttt{run\_id} 登记配置、数据版本、Git 提交号、随机种子、指标与结论（仓库 experiments/index.csv），论文中的数字均可追溯到对应运行目录。随机性来源统一固定（如 $A_1$ 拟合/留出划分使用 seed=20260923 的记录 ID 哈希规则），不存在未登记的随机优化器。

\subsection{问题一实验设置}
\textbf{质量评价}：六套方案构成敏感性矩阵——CRITIC--TOPSIS（主方案）、CRITIC 加权和、等权加权和、等权 TOPSIS、删除 DSIR 的 TOPSIS、删除三个负向指标的 TOPSIS；归一化与权重仅在 $A_1$ 拟合集（40,919 条）上估计，留出集（10,311 条）与扩展集只做前向计算。\textbf{冲突消解}：阈值在 $A_1$ 拟合集上估计为 P20/P80 并冻结，备选阈值 P10/P90 与 P25/P75 用于敏感性；补偿参数 $\lambda\in\{0,0.1,0.25,0.5\}$；输出候选冲突覆盖率、强度、广度与候选评分四张结果表。\textbf{配比建模}：A4 五折交叉验证选择 Ridge 正则强度；A6/A7 同尺度、A8/A9 与 A10/A11 跨尺度检验；A12--A15 外推诊断；基线敏感性覆盖参考点选择（参考坐标对 Helmert-ilr）与零值替换方案。

\subsection{问题二至问题四的实验状态}
问题二已完成经典 N--D 标度律基线（M0）与质量敏感性模型（M1）的探索性拟合、模型内弹性审计与 Q/p 情景包（运行号见 experiments/index.csv 的 T-Q2-MIGRATION 系列）；受数据条件限制，A、B 附件无行级关联键，正式联合拟合不成立，问题二以条件接口方式使用问题一的 $Q_g$。问题三已完成将冻结的 M0/M1 参数传播入优化器的参数不确定性情景分析（120+225 组）。问题四已完成 C8 逐任务 JSON 的截断审计与字段合同核验。上述实验的正文化呈现随相应章节定稿补充。
